% !TEX root = ../main.tex
\begin{abstract}
A policy's rank on nuScenes' open-loop metrics or NAVSIM's leaderboard score predicts its rank at a
new deployment site only as well as a coin flip ($\rho=-0.54$--$+0.49$ on near-pedestrian scenes): a
benchmark score reports an outcome, not whether a vulnerable road user had anything to do with it. We
propose a \emph{check-up} --- a few hundred verified, pixel-level counterfactual edits (removing the
pedestrian, re-rendering the scene at night) read as a diagnosis rather than a ranking. Raw response
magnitude is a poor safety signal: across \NumSGscenes{} near-pedestrian scenes and six policies, only
\NumRealPct\% of responses are genuine avoidance by a three-part test, though a near-zero response
still reliably certifies that the pedestrian was ignored (\NumPFbig\% of scenes). Moving from
left-hand to right-hand drive, what transfers is not how strongly a policy reacts but how steadily it
drives, and that exam beats the leaderboard score itself at predicting the new side's behaviour
($\rho={\NumSpBeh}$). A pre-registered, sample-size-priced transfer test confirms the operating range:
orderings and verdicts carry over, point values do not.
\end{abstract}
